% =====================================================================
% BMED365: Practical Skills
% Beamer presentation - Topic List F01-F14
% =====================================================================
\documentclass[aspectratio=169, 10pt]{beamer}

% =====================================================================
% PACKAGES
% =====================================================================
\usepackage[utf8]{inputenc}
\usepackage[T1]{fontenc}
\usepackage[english]{babel}
\usepackage{graphicx}
\usepackage{tikz}
\usetikzlibrary{shapes.geometric, arrows, positioning, calc}
\usepackage{booktabs}
\usepackage{amsmath}
\usepackage{fontawesome5}
\usepackage{listings}
\usepackage{xcolor}
\usepackage{hyperref}
\hypersetup{colorlinks=true, linkcolor=blue, urlcolor=blue}

% Code formatting
\definecolor{codebg}{rgb}{0.95,0.95,0.92}
\lstset{
    basicstyle=\ttfamily\footnotesize,
    breaklines=true,
    frame=single,
    backgroundcolor=\color{codebg},
}

% =====================================================================
% THEME AND COLORS
% =====================================================================
\usetheme{Madrid}
\usecolortheme{beaver}

% =====================================================================
% TITLE INFO
% =====================================================================
\title{Practical Skills}
\subtitle{BMED365: Topic List F01--F14}
\author{BMED365}
\date{Spring 2026}

% =====================================================================
% DOCUMENT
% =====================================================================
\begin{document}

% Title page
\begin{frame}
    \titlepage
\end{frame}

% Table of Contents
\begin{frame}{Overview}
    \begin{enumerate}
        \item \textbf{Jupyter Notebooks and Google Colab}
        (\href{https://jupyter.org/}{Jupyter}, \href{https://colab.research.google.com/}{Colab})
        \item \textbf{Python Fundamentals}
        (\href{https://numpy.org/}{NumPy}, \href{https://pandas.pydata.org/}{Pandas}, \href{https://matplotlib.org/}{Matplotlib})
        \item \textbf{Machine Learning with scikit-learn}
        (\href{https://scikit-learn.org/}{scikit-learn})
        \item \textbf{Network Analysis with NetworkX}
        (\href{https://networkx.org/}{NetworkX})
        \item \textbf{Deep Learning with PyTorch}
        (\href{https://pytorch.org/}{PyTorch})
        \item \textbf{AI Tools and LaTeX}
        (\href{https://chat.openai.com/}{ChatGPT}, \href{https://gemini.google.com/}{Gemini}, \href{https://claude.ai/}{Claude}, \href{https://www.overleaf.com/}{Overleaf})
    \end{enumerate}
\end{frame}

% =====================================================================
% SECTION: Jupyter and Colab
% =====================================================================
\section{Jupyter Notebooks and Google Colab}

\begin{frame}{F01: Run Jupyter Notebooks in Google Colab}
    \textbf{What is Jupyter Notebooks?}
    \begin{itemize}
        \item Interactive documents that combine code, text, and visualizations
        \item Run Python code cell by cell
        \item Standard in data science work
    \end{itemize}
    
    \vspace{0.3cm}
    \textbf{Google Colab:}
    \begin{itemize}
        \item Free cloud-based Jupyter environment from Google
        \item No installation - runs in the browser
        \item Free GPU access (important for deep learning!)
        \item Integrated with Google Drive
    \end{itemize}
    
    \vspace{0.3cm}
    \textbf{Getting started:}
    \begin{enumerate}
        \item Go to \texttt{colab.research.google.com}
        \item Log in with Google account
        \item ``New notebook'' or open from GitHub
    \end{enumerate}
    
    \begin{block}{\footnotesize Tip}
        \footnotesize Enable GPU: Runtime $\rightarrow$ Change runtime type $\rightarrow$ GPU
    \end{block}
\end{frame}

% =====================================================================
% SECTION: Python Fundamentals
% =====================================================================
\section{Python Fundamentals}

\begin{frame}[fragile]{F02: Use Python variables, lists, and simple functions}
    \textbf{Variables:}
    \begin{lstlisting}[language=Python]
age = 45            # int
name = "Patient A"  # str
risk = 0.73         # float
is_sick = True      # bool
    \end{lstlisting}
    
    \textbf{Lists:}
    \begin{lstlisting}[language=Python]
symptoms = ["headache", "nausea", "fatigue"]
values = [1.2, 3.4, 5.6, 7.8]
symptoms.append("fever")  # Add element
    \end{lstlisting}
    
    \textbf{Functions:}
    \begin{lstlisting}[language=Python]
def calculate_bmi(weight, height):
    return weight / (height ** 2)

bmi = calculate_bmi(70, 1.75)  # -> 22.9
    \end{lstlisting}
\end{frame}

\begin{frame}[fragile]{F03: Import and use libraries (numpy, pandas, matplotlib)}
    \textbf{Importing libraries:}
    \begin{lstlisting}[language=Python]
import numpy as np
import pandas as pd
import matplotlib.pyplot as plt
    \end{lstlisting}
    
    \vspace{0.3cm}
    \begin{columns}[T]
        \begin{column}{0.48\textwidth}
            \textbf{NumPy:}
            \begin{itemize}
                \item Numerical computations
                \item Arrays and matrices
                \item Mathematical functions
            \end{itemize}
            \begin{lstlisting}[language=Python]
arr = np.array([1, 2, 3])
mean = np.mean(arr)
            \end{lstlisting}
        \end{column}
        \begin{column}{0.48\textwidth}
            \textbf{Pandas:}
            \begin{itemize}
                \item Data manipulation
                \item DataFrames (tables)
                \item CSV, Excel I/O
            \end{itemize}
            \begin{lstlisting}[language=Python]
df = pd.read_csv("data.csv")
df.head()
            \end{lstlisting}
        \end{column}
    \end{columns}
\end{frame}

\begin{frame}[fragile]{F04: Read and inspect datasets with pandas}
    \textbf{Reading data:}
    \begin{lstlisting}[language=Python]
df = pd.read_csv("patients.csv")
    \end{lstlisting}
    
    \textbf{Inspecting data:}
    \begin{lstlisting}[language=Python]
df.head()        # First 5 rows
df.info()        # Column types, null values
df.describe()    # Statistics for numeric columns
df.shape         # (number of rows, number of columns)
df.columns       # Column names
    \end{lstlisting}
    
    \textbf{Filtering and selection:}
    \begin{lstlisting}[language=Python]
# Select columns
df[["age", "diagnosis"]]

# Filter rows
elderly = df[df["age"] > 65]
    \end{lstlisting}
\end{frame}

% =====================================================================
% SECTION: Machine Learning with scikit-learn
% =====================================================================
\section{Machine Learning with scikit-learn}

\begin{frame}[fragile]{F05: Train a simple model with scikit-learn}
    \textbf{Typical ML workflow:}
    \begin{lstlisting}[language=Python]
from sklearn.model_selection import train_test_split
from sklearn.tree import DecisionTreeClassifier
from sklearn.metrics import accuracy_score

# 1. Prepare data
X = df[["feature1", "feature2"]]
y = df["label"]

# 2. Split into train and test
X_train, X_test, y_train, y_test = train_test_split(
    X, y, test_size=0.2, random_state=42)

# 3. Create and train model
model = DecisionTreeClassifier()
model.fit(X_train, y_train)

# 4. Evaluate
y_pred = model.predict(X_test)
print(f"Accuracy: {accuracy_score(y_test, y_pred):.2f}")
    \end{lstlisting}
\end{frame}

\begin{frame}[fragile]{F06: Visualize results with matplotlib}
    \textbf{Basic plotting:}
    \begin{lstlisting}[language=Python]
import matplotlib.pyplot as plt

# Line chart
plt.plot([1, 2, 3, 4], [10, 20, 25, 30])
plt.xlabel("Time"); plt.ylabel("Value")
plt.title("My figure")
plt.show()
    \end{lstlisting}
    
    \begin{columns}[T]
        \begin{column}{0.48\textwidth}
            \textbf{Histogram:}
            \begin{lstlisting}[language=Python]
plt.hist(df["age"], bins=20)
plt.xlabel("Age")
plt.show()
            \end{lstlisting}
        \end{column}
        \begin{column}{0.48\textwidth}
            \textbf{Scatter plot:}
            \begin{lstlisting}[language=Python]
plt.scatter(df["x"], df["y"])
plt.xlabel("X"); plt.ylabel("Y")
plt.show()
            \end{lstlisting}
        \end{column}
    \end{columns}
    
    \vspace{0.2cm}
    \begin{block}{\footnotesize \href{https://seaborn.pydata.org/}{Seaborn}}
        \footnotesize \texttt{import seaborn as sns} - Beautiful visualizations with simple code!
    \end{block}
\end{frame}

% =====================================================================
% SECTION: Network Analysis
% =====================================================================
\section{Network Analysis with NetworkX}

\begin{frame}[fragile]{F07: Use \href{https://networkx.org/}{NetworkX} for simple network analysis}
    \textbf{Create and manipulate graphs:}
    \begin{lstlisting}[language=Python]
import networkx as nx

# Create graph
G = nx.Graph()
G.add_nodes_from(["P1", "P2", "P3", "P4"])
G.add_edges_from([("P1", "P2"), ("P2", "P3"), ("P1", "P3")])

# Basic analysis
print(f"Number of nodes: {G.number_of_nodes()}")
print(f"Number of edges: {G.number_of_edges()}")

# Centrality
deg_cent = nx.degree_centrality(G)
print(f"Degree centrality: {deg_cent}")

# Visualization
nx.draw(G, with_labels=True, node_color="lightblue")
plt.show()
    \end{lstlisting}
\end{frame}

% =====================================================================
% SECTION: Deep Learning
% =====================================================================
\section{Deep Learning with PyTorch}

\begin{frame}[fragile]{F08: Build and train a model with \href{https://pytorch.org/}{PyTorch}}
    \textbf{Simple MLP in PyTorch:}
    \begin{lstlisting}[language=Python,basicstyle=\ttfamily\scriptsize]
import torch
import torch.nn as nn

class SimpleMLP(nn.Module):
    def __init__(self, input_size, hidden_size, num_classes):
        super().__init__()
        self.fc1 = nn.Linear(input_size, hidden_size)
        self.relu = nn.ReLU()
        self.fc2 = nn.Linear(hidden_size, num_classes)

    def forward(self, x):
        out = self.fc1(x)
        out = self.relu(out)
        out = self.fc2(out)
        return out

model = SimpleMLP(784, 128, 10)  # MNIST example
    \end{lstlisting}

    \begin{block}{\footnotesize Lab 2}
        \footnotesize Full training with loss, optimizer, and training loop is covered in Lab 2.
    \end{block}
\end{frame}

% =====================================================================
% SECTION: AI Tools
% =====================================================================
\section{AI Tools and LaTeX}

\begin{frame}{F09: Use AI tools as coding assistants}
    \begin{columns}[T]
        \begin{column}{0.52\textwidth}
            \textbf{Useful applications:}
            \begin{itemize}
                \item \faCode~\textbf{Explain code:} ``What does this function do?''
                \item \faBug~\textbf{Debugging:} ``Why am I getting this error?''
                \item \faLightbulb~\textbf{Suggestions:} ``Make this more efficient?''
                \item \faPuzzlePiece~\textbf{Generate code:} ``Write a function...''
            \end{itemize}
        \end{column}
        \begin{column}{0.45\textwidth}
            \textbf{Tips for effective use:}
            \begin{enumerate}
                \item Be \textbf{specific} in your questions
                \item Include \textbf{context}
                \item Always \textbf{verify} AI code
                \item Use as a \textbf{learning tool}
            \end{enumerate}
        \end{column}
    \end{columns}

    \vspace{0.1cm}
    \begin{block}{\footnotesize AI-powered IDEs (Integrated Development Environments)}
        \scriptsize
        \textbf{IDE} = Editor + compilation + debugger + tools in one.
        \textbf{AI-powered IDE} = IDE with LLM agents that can plan, write, test and validate code autonomously.\\
        \textit{Significance for medical/biomedical research:} Lowers the barrier to develop analysis tools, accelerates prototyping of ML models, and enables faster translation from research to clinic.\\[3pt]
        \begin{tabular}{@{}ll@{}}
        \href{https://www.cursor.com/}{Cursor}, \href{https://windsurf.com/}{Windsurf} & Editors with agents, plan mode \\
        \href{https://antigravity.dev/}{Google Antigravity} & Agent-first, multi-model (Gemini 3) \\
        \href{https://replit.com/}{Replit Agent} & Cloud-based IDE with AI agent \\
        \href{https://lovable.dev/}{Lovable} & Natural language $\rightarrow$ full-stack app \\
        \href{https://github.com/features/copilot}{GitHub Copilot}, \href{https://www.jetbrains.com/ai/}{JetBrains AI} & AI assistants in existing IDEs
        \end{tabular}
    \end{block}
\end{frame}

\begin{frame}[fragile]{F10: Write documents with \href{https://www.latex-project.org/}{LaTeX}/\href{https://www.overleaf.com/}{Overleaf}}
    \begin{columns}[T]
        \begin{column}{0.28\textwidth}
            \textbf{\href{http://www.bibtex.org/}{BibTeX} (.bib file):}
            \begin{lstlisting}[language=TeX,basicstyle=\ttfamily\tiny]
@article{smith2024,
  author  = {Smith, J.},
  title   = {AI in Medicine},
  journal = {Nature Med.},
  year    = {2024},
  volume  = {30},
  pages   = {123--130},
  doi     = {10.1038/s41591-
             024-01234-5}
}

@book{bishop2006,
  author    = {Bishop, C.},
  title     = {Pattern Recog.},
  publisher = {Springer},
  year      = {2006},
  isbn      = {978-0387310732}
}
            \end{lstlisting}
            \tiny In text: \texttt{\textbackslash cite\{smith2024\}}
        \end{column}
        \begin{column}{0.40\textwidth}
            \textbf{IMRAD structure:}
            \begin{lstlisting}[language=TeX,basicstyle=\ttfamily\tiny,literate={å}{{\aa}}1]
\documentclass{article}
\usepackage{graphicx,booktabs}
\begin{document}
\title{Title}
\author{Author}
\maketitle
\begin{abstract} ... \end{abstract}
% Introduction
\section{Introduction}
Background and aim \cite{smith2024}.
% Materials and methods
\section{Methods}
Data collection and analysis.
% Results
\section{Results}
\begin{table}[h]
\begin{tabular}{lcc} \toprule
Model & AUC & F1 \\ \midrule
CNN & 0.92 & 0.88 \\ \bottomrule
\end{tabular}
\caption{Results}\label{tab:1}
\end{table}
% Discussion
\section{Discussion}
Implications and limitations.
\bibliography{refs}
\end{document}
            \end{lstlisting}
        \end{column}
        \begin{column}{0.30\textwidth}
            \textbf{Journals with LaTeX:}
            \tiny
            \begin{itemize}
                \item \href{https://journals.plos.org/plosmedicine/s/latex}{PLOS Medicine}
                \item \href{https://www.nature.com/nature/for-authors/formatting-guide}{Nature / Nat. Med.}
                \item \href{https://academic.oup.com/bioinformatics/pages/instructions_for_authors}{Bioinformatics}
                \item \href{https://www.thelancet.com/pb/assets/raw/Lancet/authors/tl-info-for-authors.pdf}{The Lancet}
                \item \href{https://www.bmj.com/about-bmj/resources-authors}{BMJ}
                \item \href{https://www.elsevier.com/researcher/author/policies-and-guidelines/latex-instructions}{Elsevier (many)}
                \item \href{https://www.springer.com/gp/livingreviews/latex-templates}{Springer Nature}
                \item \href{https://www.mdpi.com/authors/latex}{MDPI (open access)}
                \item \href{https://ieee-dataport.org/sites/default/files/analysis/27/IEEE-Reference-Guide.pdf}{IEEE JBHI}
                \item \href{https://www.jmlr.org/format/format.html}{JMLR}
                \item \href{https://www.frontiersin.org/guidelines/author-guidelines}{Frontiers}
            \end{itemize}
            \vspace{0.05cm}
            \scriptsize \href{https://www.overleaf.com/gallery/tagged/academic-journal}{Overleaf Gallery}
        \end{column}
    \end{columns}
\end{frame}

% =====================================================================
% SECTION: Context Engineering
% =====================================================================
\section{Context Engineering for LLM}

\begin{frame}{F11: Compose effective prompts for medical tasks}
    \textbf{From Prompt Engineering to Context Engineering:}
    \begin{itemize}
        \item \textbf{Prompt engineering:} Crafting good questions and instructions
        \item \textbf{Context engineering:} Designing the \textit{entire context} the model uses
    \end{itemize}
    
    \vspace{0.2cm}
    \textbf{Components of context engineering:}
    \begin{enumerate}
        \item \textbf{System prompt:} Role and basic instructions
        \item \textbf{User prompt:} The specific question
        \item \textbf{Retrieved context (RAG):} Medical records, guidelines, articles
        \item \textbf{Tools:} Lab lookup, drug database, ICD-10
        \item \textbf{Conversation history:} Previous dialog in the consultation
        \item \textbf{Structured data:} JSON with lab values, vital signs
    \end{enumerate}
    
    \begin{block}{\footnotesize Medical significance}
        \footnotesize Rich context = more relevant, personalized, and safe AI responses
    \end{block}
\end{frame}

\begin{frame}{F12: Apply context engineering in practice}
    \textbf{Example: Simple prompt vs. rich context}
    
    \vspace{0.2cm}
    \begin{columns}[T]
        \begin{column}{0.45\textwidth}
            \begin{alertblock}{\footnotesize Simple prompt}
                \scriptsize
                \texttt{What is the treatment for heart failure?}
                
                \vspace{0.1cm}
                $\rightarrow$ Generic answer, not personalized
            \end{alertblock}
        \end{column}
        \begin{column}{0.52\textwidth}
            \begin{block}{\footnotesize Rich context}
                \scriptsize
                \texttt{System: Clinical decision support} \\
                \texttt{Patient: 72 years, eGFR 45, Warfarin} \\
                \texttt{RAG: ESC Guidelines 2023} \\
                \texttt{Question: Treatment changes?}
                
                \vspace{0.1cm}
                $\rightarrow$ Personalized, evidence-based answer
            \end{block}
        \end{column}
    \end{columns}
    
    \vspace{0.2cm}
    \textbf{Skill levels:}
    \footnotesize
    \begin{tabular}{ll}
        \textbf{Beginner:} & Prompt formulation (role, format, examples) \\
        \textbf{Intermediate:} & Structured system prompts with sections \\
        \textbf{Advanced:} & RAG integration and tool use \\
        \textbf{Expert:} & Multi-agent orchestration \\
    \end{tabular}
\end{frame}

\begin{frame}{F13: Choose the right LLM model for the task}
    \textbf{Model strengths for medical tasks:}
    
    \vspace{0.2cm}
    \footnotesize
    \begin{tabular}{lll}
        \toprule
        \textbf{Task type} & \textbf{Recommended model} & \textbf{Rationale} \\
        \midrule
        Structured extraction & GPT-5.2 & Good instruction following \\
        Ethical dilemmas & Claude Opus 4.5 & Nuanced reasoning \\
        Image analysis & Gemini 3 & Multimodal from the ground up \\
        Long record analysis & Claude Opus 4.5 & 200K token context \\
        Updated research & Gemini 3 & Integrated real-time search \\
        Tool integration & GPT-5.2 & Good API for functions \\
        Patient communication & Claude Opus 4.5 & Empathetic, patient-centered \\
        Calculations & Gemini 3 & Strong in science \\
        \bottomrule
    \end{tabular}
    
    \vspace{0.2cm}
    \normalsize
    \begin{block}{\footnotesize Practical recommendation}
        \footnotesize For critical tasks: Test on multiple models and compare output
    \end{block}
\end{frame}

\begin{frame}{F14: Compare output from different LLMs}
    \textbf{Evaluation criteria for medical LLM output:}
    
    \vspace{0.2cm}
    \begin{columns}[T]
        \begin{column}{0.48\textwidth}
            \textbf{Quality dimensions:}
            \begin{enumerate}
                \item \textbf{Medical accuracy} \\
                      {\scriptsize Is the information correct?}
                \item \textbf{Uncertainty handling} \\
                      {\scriptsize Is uncertainty expressed appropriately?}
                \item \textbf{Structured output} \\
                      {\scriptsize Are format instructions followed?}
                \item \textbf{Safety behavior} \\
                      {\scriptsize Is referral to physician recommended when needed?}
                \item \textbf{Clinical utility} \\
                      {\scriptsize Is the answer practically applicable?}
            \end{enumerate}
        \end{column}
        \begin{column}{0.48\textwidth}
            \textbf{Practical approach:}
            \begin{enumerate}
                \item Define test case with ground truth
                \item Send same prompt to 2--3 models
                \item Evaluate against criteria (1--5 scale)
                \item Document strengths/weaknesses
                \item Choose model for production
            \end{enumerate}
            
            \vspace{0.2cm}
            \begin{alertblock}{\footnotesize Lab 3 Notebook 04}
                \scriptsize Contains ready-made experiments for model comparison
            \end{alertblock}
        \end{column}
    \end{columns}
\end{frame}

% =====================================================================
% SUMMARY
% =====================================================================
\section*{Summary}

\begin{frame}{Summary: Practical Skills}
    \textbf{Jupyter and Python:}
    \begin{itemize}
        \item F01--F04: Colab, variables, lists, functions, pandas
    \end{itemize}
    
    \textbf{Machine Learning:}
    \begin{itemize}
        \item F05--F06: scikit-learn workflow, matplotlib visualization
    \end{itemize}
    
    \textbf{Specialized:}
    \begin{itemize}
        \item F07: NetworkX for network analysis
        \item F08: PyTorch for deep learning
    \end{itemize}
    
    \textbf{Tools and LLM:}
    \begin{itemize}
        \item F09--F10: AI assistants, LaTeX/Overleaf
        \item F11--F14: Context engineering, model selection, prompt optimization
    \end{itemize}
    
    \vspace{0.2cm}
    \begin{block}{QuickStart and Labs}
        All these skills are practiced through the QuickStart and Labs 0--3. Practice makes perfect!
    \end{block}
\end{frame}

\end{document}
